\section{Introduction}

The heat kernel on a compact Lie group is simultaneously a spectral object, a
short-time Gaussian, and a convolution kernel.  These descriptions expose
different structures.  Character sums make normalization and the semigroup
law transparent; Poisson images expose the ultraviolet Gaussian; neither
description by itself is a stable formula at every class angle and heat time.
The negative logarithm compounds the problem: differentiating it requires a
strict lower bound before quotient estimates become legitimate.

This issue is concrete in lattice gauge theory.  The heat-kernel action of
Menotti and Onofri is classical \cite{MenottiOnofri1981}, as are its relation
to the Manton action \cite{Manton1980} and its exact convolution behavior in
two-dimensional Yang--Mills theory \cite{GrossKingSengupta1989,Sengupta1993}.
Constructive and validated estimates require more specific data: a declared
normalization, a closed interval, derivatives of $-\log p_t$ through a
prescribed order, an explicit truncation rule, and a rigorously controlled
positive lower bound.

The central result of this note is a global version of these estimates for
$\SU$.  Given any finite order $M$, the theorem produces explicit bounds on
three overlapping class-angle charts.  The identity and transition charts
use the same paired correction with different time majorants.  The cut chart
uses positive odd-image pairs and a principal-pair lower bound.  The
construction is not a new formula for the heat kernel.  Its contribution is
an explicit finite-time enclosure calculus, including the antipodal
Morse--Bott prefactor and an exact-rational numerical realization.

The proof has four ingredients.
\begin{enumerate}[label=(\roman*)]
\item Opposite images are paired before division by the class angle.
\item An integral identity for $\sinh u/u$ gives a derivative majorant with
  one expression valid at the identity and in the interior.
\item Exponential partial Bell polynomials transfer the hierarchy through the
  logarithm without repeating an order-by-order quotient calculation.
\item At the antipode, the two shortest radial image terms are paired before
  taking a removable limit; the resulting prefactor reflects the continuous
  family of minimizing geodesics.
\end{enumerate}

The paper is organized from exact formulas to validated consequences.
Sections~\ref{sec:conventions}--\ref{sec:arbitrary} fix conventions and prove
the derivative hierarchy.  Sections~\ref{sec:action} and
\ref{sec:endpoints} treat the action Hessian and endpoint charts.
Section~\ref{sec:tails} proves normalized tails.  Section~\ref{sec:certificate}
describes the rational validation, and Section~\ref{sec:applications}
derives the lattice normalization and a concrete plaquette estimate.
Complete summation, endpoint, and validation details are deferred to the
appendices.
